\documentclass[11pt,fleqn]{report}
\usepackage{graphicx}
\usepackage{amssymb}
\usepackage{tabularx}
\usepackage[T1]{fontenc}
\usepackage{enumitem}
\usepackage{makecell}
\usepackage{comment}
\usepackage{float}
\usepackage{pdflscape}
\usepackage{adjustbox}
\usepackage{ragged2e}
% \usepackage{tabularray}
\usepackage[skins]{tcolorbox}
\usepackage{lmodern}
\usepackage{cml}

\lstnewenvironment{codeXML}
    {\lstset{basicstyle=\small\ttfamily,
             language=XML,
             xleftmargin=0.0in,
             xrightmargin=0.0in,
             backgroundcolor=\color{Orchid!60}, % 60% orchid, 40% white
             escapeinside={\%*}{*)}, %  .tex file characters inside %*...*) are
                                     %  treated as Tex commands.
            }
    }
    { }

\newcommand{\ModelName}{CML Two-Level Vent}
\newcommand{\ModelAuthor}{Suyog Benegalrao \& Alexandre Masset}
\newcommand{\ModelAffiliation}{NASA}
\newcommand{\ModelDate}{September 2025}
\pagestyle{plain}
\usepackage[parfill]{parskip}
\usepackage[margin=1.5cm]{caption}
\usepackage{subfig}

\setlist[enumerate, 1]{label=\arabic*., ref=\arabic*}
\setlist[enumerate, 2]{label=\arabic{enumi}.\arabic*., 
                       ref=\arabic{enumi}.\arabic*}
\setlist[enumerate, 3]{label=\arabic{enumi}.\arabic{enumii}.\arabic*., 
                       ref=\arabic{enumi}.\arabic{enumii}.\arabic*}

\begin{document}
\pagestyle{empty}
\titlePage
\pagenumbering{roman}
\tableofcontents % Print table of contents
\vfill
\pagebreak
{\Large \textbf {Revision History}}
{\large
\begin{center}
\begin{tabular}{|l||l|l|l|}
\hline
 \textbf{Version} & \textbf{Date} & \textbf{Author} & \textbf{Purpose} \\ \hline \hline
1 & Sep 2025 & Suyog Benegalrao \& Alexandre Masset & Initial Version \\ \hline
\end{tabular}
\end{center}
}

\chapter{Introduction} % Make a "chapter" heading
\pagenumbering{arabic} % Start text with arabic 1
\noindent 
This model is an extension of the CML Vent model, with the specific purpose to model venting events
of a candidate vehicle. Indeed, this model was built to answer RPOD requirements but can be used for any
vehicle if the user finds the proposed behavior appropriate for their application.

More specifically, this model proposes :
\begin{itemize}
    \item A SimpleVent class extension with two venting modes: Low and High.
    \item Parameters for appropriate Monte Carlo simulation setup
    \item A VentSet class extension to manage these venting modes
\end{itemize}

The following documentation assumes that the user is familiar with the parent CML Vent model and shall
refer to its documentation if needed as this model explicitly overrides and uses the parent's methods.

\chapter{Requirements}
\noindent The requirements for the two-level venting model are:

\begin{enumerate}[noitemsep]
    \item At the model-level, the CML Two-Level Vent model shall:
    \begin{enumerate}
        \item meet the CML Vent parent model requirements.
        \item provide implementation of individual two-level vents and systems of two-level vents.
        \item provision each two-level vent with a set of boolean flags providing:
        \begin{enumerate}
            \item a user-interface for configuring and verifying the configuration and operation of the two-level vent.
            \item an information-only indication whether the two-level vent is operating in low mode or high mode.
        \end{enumerate}
        \item provide appropriate feedback and error-catching when configurations are incorrect
    \end{enumerate}
    \item A two-level vent-set shall:
    \begin{enumerate}
        \item be associated with a vehicle
        \item be associated with a number of vents
        \item provide access to the two-level vent instances it controls
        \item handle start/stop of the vents it controls
        \item strictly contain TwoLevelVent instances
    \end{enumerate}
        \item A two-level vent shall:
        \begin{enumerate}
            \item be associated with a single two-level vent-set
            \item be configured as dynamic to ensure integration of forces on the vehicle body over time
            \item have a low venting mode, enforcing a user-defined minimum acceleration
            \item have a high venting mode, enforcing a delta-V based on user-defined acceleration and high-mode event durations
            \item automate transition from a venting mode to another 
            \item transition to a low-mode when:
            \begin{enumerate}
                \item the current high-mode firing duration is reached
                \item the user-defined maximum delta-V value is reached for protection
                \item the vent stops venting
                \item the vent is reset
                \item the vent is disabled
                \item the two-level vent-set is deactivated through unsubscribtion
                \item the high-mode duration is undefined or zero
            \end{enumerate}
            \item transition to a high-mode when:
            \begin{enumerate}
                \item the two-level vent starts venting.
                \item the current high-mode wait time (period) is reached
            \end{enumerate}
            \item be configurable with parameters describing:
            \begin{enumerate}
                \item the period between the start of high mode events
                \item the duration of high mode venting events
                \item the minimum acceleration to use in low-venting mode
                \item the maximum delta-V the vehicle cannot exceed in one integration step during high-venting mode
                \item the acceleration applied to the vehicle during a high-venting event
            \end{enumerate}
            \item provide a set of parameters for Monte Carlo simulations such as:
            \begin{enumerate}
                \item minimum and maximum wait time between start of high-mode venting events
                \item minimum and maximum duration of high-mode venting events
            \end{enumerate}
            \item provide flags to enable/disable randomization of the MC parameters
            \item provide an interface to reset scheduled high-mode events and transition to a low venting mode in a single event
        \end{enumerate}
\end{enumerate}

% \makecell{Two\_Level\_Venting\_Model-0001 \\ Vent Number} &
% \underline{Requirement}: The Two-Level Venting Model shall model multiple, independent active venting disturbance effects. \newline
% \underline{Rationale}: Two-Level Vents can potentially produce periodic venting effects due to various effects. This will perturb the vehicle's dynamic behavior by applying a disturbance acceleration. The number and cadence of these active venting events is unknown, so the Two-Level Venting Model will include a configurable static value representing the maximum number of vents, and each vent active disturbance is modeled independently such that multiple, independent vent active disturbances can occur at any cadence. \\

% \makecell{Two\_Level\_Venting\_Model-0002 \\ Active Venting Disturbance Direction} &
% \underline{Requirement}: The Two-Level Venting Model shall prescribe a configurable, constant unit vector representing the direction of the perturbation caused by each independent active venting disturbance. \newline
% \underline{Rationale}: When an active venting event occurs in the Two-Level Venting Model, the direction of the disturbance acceleration will be defined as a constant unit vector for each independent vent.  The unit vector origin (tail) will be placed at the Two-Level center-of-gravity.   \\

% \makecell{Two\_Level\_Venting\_Model-0003 \\ Active Venting Disturbance Acceleration Magnitude} &
% \underline{Requirement}: The Two-Level Venting Model shall include a configurable, constant parameter representing the acceleration magnitude of the perturbation caused by each independent active venting disturbance. \newline
% \underline{Rationale}: When an active venting event occurs in the Two-Level Venting Model, the magnitude of the disturbance will be defined by a constant parameter for each independent vent.  The acceleration disturbance vector is calculated by multiplying the acceleration magnitude by the pre-defined unit vector. \\

% \makecell{Two\_Level\_Venting\_Model-0004 \\ Active Venting Disturbance Duration} &
% \underline{Requirement}: The Two-Level Venting Model shall model the disturbance duration for each active venting event independently such that it is either randomly sampled from a range of configurable maximum and minimum duration values or set to a pre-defined configurable fixed value. \newline
% \underline{Rationale}: For each active venting event, the disturbance duration is either randomized or fixed.  The resultant Delta-V effect is achieved through simulation integration. \\

% \makecell{Two\_Level\_Venting\_Model-0005 \\ Active Venting Disturbance Period} &
% \underline{Requirement}: The Two-Level Venting Model shall institute a given wait-time between each independent active venting disturbance event that is either randomly sampled from a range of configurable maximum and minimum wait-time values or set to a fixed value. \newline
% \underline{Rationale}: For each active venting event, the next active venting event is scheduled such that it occurs at a either a randomly sampled time or a fixed value. The wait-time / period is defined as the time between the start times of corresponding active venting events. \\


\chapter{Model Architecture}
\noindent Most of the model architecture is inherited from the Common Model Libraries (CML) VentSet
and SimpleVent. This section covers the members and functions that are unique to the Two-Level Vent model.

\section{TwoLevelVent Class}
\subsection{TwoLevelVent Class Members}
\noindent The tables below contain the members for the TwoLevelVent class. This class inherits from the 
SimpleVent class.

\begin{table}[!h]
    \caption{TwoLevelVent members for Monte Carlo setup}
    \begin{tabularx}{\textwidth}{|c|c|X|X|} 
    \hline
    \textbf{Access} & \textbf{Type} & \textbf{Name} & \textbf{Purpose} \\
    \hline
    public & double & wait\_time\_min & Minimum wait time between start of high mode firings \\
    \hline
    public & double & wait\_time\_max & Maximum wait time between start of high mode firings \\
    \hline
    public & double & high\_duration\_min & Minimum duration of high mode firings \\
    \hline
    public & double & high\_duration\_max & Maximum duration of high mode firings \\
    \hline
    public & bool & wait\_time\_randFlag & Flag where TRUE randomizes at each update the wait time between
                                        the specified min and max values \\
    \hline
    public & bool & high\_duration\_randFlag & Flag where TRUE randomizes at each update the duration between
                                        the specified min and max values \\
    \hline
    \end{tabularx}
\end{table}

\begin{table}[!h]
    \caption{Other TwoLevelVent members}
    \begin{tabularx}{\textwidth}{|c|c|X|X|} 
    \hline
    \textbf{Access} & \textbf{Type} & \textbf{Name} & \textbf{Purpose} \\
    \hline
    public & double & accel\_mag\_high & Magnitude of high mode firing acceleration \\
    \hline
    public & double & accel\_mag\_low & Magnitude of low mode firing acceleration \\
    \hline
    public & double & high\_duration & Duration of high mode firings. Updated if high\_duration\_randFlag is TRUE \\
    \hline
    public & double & high\_wait\_time & Wait time of high mode firings. Updated if wait\_time\_randFlag is TRUE \\
    \hline
    public & double & dv\_max & Maximum delta-V allowed for each high mode firing \\
    \hline
    public & TwoLevelVentReset & action\_reset & Event-action instance to reset vent configuration to low-mode on call \\
    \hline
    protected & const double \& & vehicle\_mass & (kg) Reference to vehicle mass \\
    \hline
    protected & std::mt19937 \& & gen & Reference to random number generator \\
    \hline
    protected & double & dv\_mag & Logged delta-V for each high mode firing \\
    \hline
    protected & double & sched\_onTime & Logged boolean whether the vent is working in high mode (1) or low-mode(0) \\
    \hline
    protected & bool & high\_venting & Logged scheduled time at which a high mode firing shall start \\
    \hline
    protected & double & start\_time\_high & Time when the current high mode firing started. \\
    \hline
    \end{tabularx}
\end{table}


\subsection{TwoLevelVent Class Methods}

\noindent The table below contains the methods for the TwoLevelVent Class.

\begin{table}[!h]
    \begin{tabularx}{\textwidth}{|c|c|X|c|X|} \hline
    \textbf{Access} & \textbf{Return Type} & \textbf{Name} & \textbf{Arguments} & \textbf{Purpose} \\
    \hline
    protected & void & check\_configuration & void & Override of SimpleVent method to check if configuration is correct with additional checks.  \\
    \hline
    protected & void & reset & void & Resets the two-level vent into a low-mode configuration, resetting the current sched\_onTime value to zero. \\
    \hline
    protected & void & stop\_venting & void & Override method that stops the vent and resets it  \\
    \hline
    protected & void & update\_force & void & Override method to compute vent force and switch between high and low mode if conditions are met.  \\
    \hline
    protected & void & high\_venting\_mode & void & Configures the vent in a high mode  \\
    \hline
    protected & void & low\_venting\_ode & void & Configures the vent in a low mode  \\
    \hline
    \end{tabularx}
\end{table}

This model also has a TwoLevelVentReset class, which is an extension of the CML EventAction class. 
More details about its use are provided in the next sections.

\section{TwoLevelVentSet Class}
\subsection{TwoLevelVentSet Class Members}
\noindent The table below contains the members for the TwoLevelVentSet class.  This class inherits from the VentSet class in CML.

\begin{table}[!h]
    \caption{TwoLevelVentSet members}
    \begin{tabularx}{\textwidth}{|c|c|X|X|} \hline
    \textbf{Access} & \textbf{Type} & \textbf{Name} & \textbf{Purpose} \\
    \hline
    public & int & rand\_seed & Random seed used for vent duration and wait time randomization \\
    \hline
    protected & const double\& & vehicle\_mass & Reference to vehicle mass (kg) \\
    \hline
    protected & std::mt19937 & gen & Random number generator for vent duration and wait time randomization \\
    \hline
    \end{tabularx}
\end{table}

\subsection{TwoLevelVentSet Class Methods}
\noindent The table below contains the methods for the TwoLevelVent Class.

\begin{table}[!h]
    \caption{TwoLevelVentSet Class Methods}
    \begin{tabularx}{\textwidth}{|c|c|X|c|X|} \hline
    \textbf{Access} & \textbf{Return Type} & \textbf{Name} & \textbf{Arguments} & \textbf{Purpose} \\
    \hline
    public & void & initialize & void & Initializes the vents and runs sanity checks on inputs  \\
    \hline
    public & TwoLevelVent * & get\_two\_level\_vent & size\_t ii & Returns pointer to the two-level vent stored at index ii  \\
    \hline
    protected & void & deactivate & void & Override method to deactivate the vent set. Called when unsubscribed. This override 
    resets all the vents to low mode so that they can start again in high mode when re-activated.\\
    \hline
    \end{tabularx}
\end{table}

\section{Mathematical Formulation}

\noindent This section will not detail how the vents compute forces. The user shall refer to the CML Vent documentation 
for more details. Instead, this section is used to explain how the two-level vent differs from its CML parent. 

\noindent A two-level vent has two venting mode, low and high.
\begin{itemize}
    \item The "low" mode corresponds to the minimum acceleration disturbance whose magnitude is flowed down from the Candidate 
    Vehicle Delta-V requirements.
    \item The "high" mode corresponds to the periodic increases in force that are meant to represent actual on-orbit venting 
    instances.
\end{itemize}

Note that in either mode, the vent is actively venting, setting its \emph{venting} parameter to true. Therefore, the \emph{duration}
parameter is set zero and the flag \emph{infinite\_duration} to true. A high venting mode shall be used when a delta-V requirement 
exists.

\noindent By default, a two-level vent is in low-mode, checking whether it shall transition to a high mode or not. The two-level vent
is meant to be dynamic, hence producing a force that is collected by the vent set. The magnitude of this force is determined by
the current vent mode.

\begin{tcolorbox}[colback=lightgray, boxrule=0pt, sharp corners]
\textbf{NOTE:} If a user specifies a venting duration, the two-level vent will interrupt any venting events
    when the duration is reached. 
\end{tcolorbox}

\begin{figure}[!h]
\begin{center}
\includegraphics[scale=0.5]{Images/example.png}
\caption{Example of expected behavior and parameterization}
\label{fig:example.png}
\end{center}
\end{figure}

\chapter{User's Guide}
\section{Instantiation and Population}

Just like its CML parent, two-level vents are meant to be controlled by a vent set. Although it is possible to populate multiple VentSets
with the same vent, TwoLevelVent instances are added at construction to a single vent set and should not be manually added to 
another vent set. Note that this model has its own override to the \emph{add\_vent()} method, set to protected here for that purpose.

\noindent A two-level vent instance requires dynamic time and the reference to the vent set it is added to in order to be initialized.
A TwoLevelVentSet instance requires a dynamic body reference and the dynamic time as well. An example is provided below.

\begin{codeC}
##include "cml/models/fhw/effectors/two_level_vent/include/two_level_vent.hh"
##include "cml/models/fhw/effectors/two_level_vent/include/two_level_vent_set.hh"
...
DynBody vehicle_body;
MassBodyInit mass_init;
TwoLevelVentSet vent_set(vehicle_body, dyn_time);	
TwoLevelVent two_level_vent(dyn_time, vent_set);	
\end{codeC}

There are two ways to create and populate a TwoLevelVentSet instance through two different constructors:
\begin{itemize}
    \item Creating independent TwoLevelVent instances that will be added at construction to the specified vent set.
    \item Specifying at construction the number of two-level vents to add to the vent set
\end{itemize}

The first method is presented above. The second method is done through:

\begin{codeC}
    size_t num_vents = 5;
    TwoLevelVentSet vent_set(num_vents, vehicle_body, dyn_time);	
\end{codeC}

\section{Configuring Vents}

If the vents are declared and instantiated outside of the TwoLevelVentSet, a user can directly modify their public parameters:

\begin{codeC}
    vent1.high_duration = 10.0;
\end{codeC}

An astute reader will notice that the constructors available are similar to those of CML VentSet class. However, it is important to note 
that the two-level vents created in the second method won't be accessible for logging as Trick cannot log STL-vectors. In this case, any vent 
configuration shall be done using the \emph{get\_two\_level\_vent(size\_t ii)} method:

\begin{codeC}
    two_level_vent.get_two_level_vent(0).high_duration = 10.0;
\end{codeC}

This also works in a Python input file. Note that the CML VentSet class also has a \emph{get\_vent(size\_t ii)} method. The main 
difference is that the latter return a SimpleVent pointer that does not have access to the TwoLevelVent parameters.

\begin{tcolorbox}[colback=lightgray, boxrule=0pt, sharp corners]
\textbf{NOTE:} The two-level vents are meant to be dynamic. The user shall not use the CML Vent methods to set the vents in an 
impulse mode. This would result in unexpected/undesired resuts.
\end{tcolorbox}

\subsection{Randomization}
The TwoLevelVent model provides the option to randomize the high mode venting firings duration and period. To do so, the user must set 
the following flag to true:

\begin{codeC}
    two_level_vent.get_two_level_vent(0).high_duration_randFlag = true;
    two_level_vent.get_two_level_vent(0).wait_time_randFLag = true;
\end{codeC}

\subsection{Event-action \& Reset}
\noindent It can be desired to reset (return to low mode and set \emph{sched\_onTime} to zero) a given vent during a run. 
The model has an event-action child class named TwoLevelVentReset that can be used for that purpose.
It is recommended to be familiar with the CML EventManager model. 

\noindent This event-action is to be added to the desired WatchValue that will trigger the TwoLevelVent \emph{reset()} method 
when its condition is met.

\noindent The event-action is stored in each two-level vent instance as the parameter \emph{action\_reset}. 
The latter must be added to the watch value.

\begin{codeC}
    WatchValuesBase<double> watch_value;
    watch_value.add_action(vent.action_reset)
\end{codeC}

\section{Initialization}
\noindent Each two-level vent initialization is handled through the two-level vent set instance. Therefore, it is only required to
run the following method to initialize all the vents added to the vent set.

\begin{codeC}
P_DYN("initialization") two_level_vent_set.initialize();	
\end{codeC}

\section{Execution}
\noindent Each two-level vent update is handled through the two-level vent set instance. Therefore, it is only required to run the 
following method to update all the vents added to the vent set.
\begin{codeC}
(DYNAMICS,"scheduled") two_level_vent_set.update();	
\end{codeC}

\section{Start/Stop}
\noindent There exists multiple options to start and stop vents. Some of them are already presented in the CML Vent model 
documentation but will be provided here for clarity.

\noindent The vent set manages all the vents, which should not be stopped/started manually through their own methods.
The vent set class has its own method for this purpose.

\noindent To start and stop all the vents, the user can call:
\begin{codeC}
    two_level_vent_set.start_vents();
    two_level_vent_set.stop_vents();
\end{codeC}

\begin{tcolorbox}[colback=lightgray, boxrule=0pt, sharp corners]
\textbf{NOTE:} These methods are defined in the CML VentSet class and call \emph{start\_venting(), stop\_venting()} methods
defined in the CML SimpleVent class. Only the latter has been overriden to reset the vent at the same time, guaranteeing that 
it starts again in a high mode when called if possible. Therefore, \emph{stop\_vents()} will call the TwoLevelVent override.
\end{tcolorbox}

\noindent To start and stop a single vent, the user can also specify an index and call:
\begin{codeC}
    two_level_vent_set.start_vent(0);
    two_level_vent_set.stop_vent(0);
\end{codeC}

\begin{tcolorbox}[colback=lightgray, boxrule=0pt, sharp corners]
\textbf{NOTE:} Same comment as above
\end{tcolorbox}

\noindent To start and stop a single vent, the user can also manually deactivate it:
\begin{codeC}
    two_level_vent_set.get_two_level_vent(0).active = false;
    two_level_vent_set.get_two_level_vent(0).active = true;
    two_level_vent_set.start_vent(0);
\end{codeC}

\begin{tcolorbox}[colback=lightgray, boxrule=0pt, sharp corners]
\textbf{NOTE:} This method is not recommended. The reason is that setting the TwoLevelVent \emph{active} parameter to false 
and back to true does not restart the vent on its own. It is required to call \emph{start\_vent(0)} method to do so. 
\end{tcolorbox}

\section{Notes on logging}
This section aims at removing confusion that arises when logging vent parameters through different tools. \\

First, note that Trick logging does not exactly match the sim dynamic time. Logging happens at the end of every Trick frame window
which has in the middle of it dynamic time incrementation. This means that the Trick frame $t=t_{frame}$  starts when the dynamic time 
$t=t_{frame} - \delta t_{step}$ is still on.

If using \emph{start\_vents()} method with \emph{trick.add\_read( i, ...)}, the update is actually applied at $t=t_{frame} - \delta t_{step}$.
This can result in shorter venting than expected since the stop time will be $t = t_{frame} - \delta t_{step} + t_{duration}$. However, the 
higher rate the sim uses, the smaller the error is, becoming totally negligible at 200Hz (most likely to be used for vents). 
Note that this issue is true for any type of vents currently implemented. \\

A second "issue" can arise from the logging tool used. One can see in the Verification part that when the vent duration is set to 2 seconds,
and the logging rate to 1Hz, there are two data points that will show high venting mode on. Koviz (and most plotting tools) will interpolate 
between the two data points before and after the vent switches to low mode. This results in a confusing interpretation that the venting is 
switching to low mode between those steps, resulting in a shorter high mode (or longer low mode). However, this is a plotting issue rather than a 
missed integration step. This flag is not updated before the next integration step is reached. Again, note that with higher logging rates, this
becomes invisible on a plot.


\chapter{Verification}

\section{Code Coverage}

\begin{figure}[!h]
\begin{center}
\includegraphics[scale=0.5]{Images/code_coverage.png}
\caption{Code Coverage}
\label{fig:code_coverage.png}
\end{center}
\end{figure}

\section{Exceptions}

\begin{itemize}
    \item In two\_level\_vent.cc, \emph{stop\_venting()}, line 131: Simple override. \emph{reset()} is tested elsewhere.
    \emph{SimpleVent::stop\_venting()} verified in CML Vent. \\
    \begin{minipage}[t]{\linewidth}
          \raggedright
          \adjustbox{valign=t}{%
            \includegraphics[scale=0.5]{Images/exception_1.png}
          }
            \captionof{figure}{two\_level\_vent\_set.cc : line 81}
          \medskip
    \end{minipage}

    \item In two\_level\_vent\_set.cc, \emph{deactivate()}, line 61: override method: methods already tested in CML Vent. \\
    \begin{minipage}[t]{\linewidth}
          \raggedright
          \adjustbox{valign=t}{%
            \includegraphics[scale=0.5]{Images/exception_2.png}
          }
            \captionof{figure}{two\_level\_vent\_set.cc : line 104}
          \medskip
    \end{minipage}

    \item In \emph{initialize()}, line 91: if the vehicle mass is zero, returns an error. Manually tested. \\
    \begin{minipage}[t]{\linewidth}
          \raggedright
          \adjustbox{valign=t}{%
            \includegraphics[scale=0.5]{Images/exception_3.png}
          }
            \captionof{figure}{two\_level\_vent\_set.cc : line 119}
          \medskip
    \end{minipage}
\end{itemize}

\newpage
\section{Unit Test Inventory}
RUN\_01\_nominal

RUN\_02\_fixedDur

RUN\_03\_fixedWaitTime

RUN\_04\_reset\_event\_action

ERROR\_rand\_config

ERROR\_rand\_duration

ERROR\_rand\_wait\_time

ERROR\_time\_config

WARN\_01\_dur

WARN\_02\_wait\_time

WARN\_03\_maxDV

FAIL\_invalid\_index


\newpage
\section{Default Unit Test Setup}

\noindent The table below contains the default parameter settings for the unit test campaign of the two-level venting model.
A single vehicle body is added at the zero position. \\

We define two vent sets with the following properties:

\subsubsection{VentSet1}
VentSet1 contains two TwoLevelVents: Vent1 and Vent2.

\begin{table}[!ht]
    \centering
    \begin{tabular}{|c|c|c|} \hline
    \textbf{Parameter Name} & \textbf{Parameter Setting} & \textbf{Units} \\
    \hline
    composite\_properties.mass & 100.0 & kg \\
    \hline
    vent1.accel\_mag\_high & 2.0 & m/s$^{2}$ \\
    \hline
    vent1.accel\_mag\_low & 1.0 & m/s$^{2}$ \\
    \hline
    vent1.high\_duration\_randFlag & TRUE & bool \\
    \hline
    vent1.high\_duration & 2.0 & s \\
    \hline
    vent1.wait\_time\_randFlag & TRUE & bool \\
    \hline
    vent1.high\_wait\_time & 5.0 & s \\
    \hline
    vent1.dv\_max & 4.0 & m/s \\
    \hline
    vent2.accel\_mag\_high & 3.0 & m/s$^{2}$ \\
    \hline
    vent2.accel\_mag\_low & 1.0 & m/s$^{2}$ \\
    \hline
    vent2.high\_duration\_randFlag & FALSE & bool \\
    \hline
    vent2.high\_duration & 2.0 & s \\
    \hline
    vent2.wait\_time\_randFlag & FALSE & bool \\
    \hline
    vent2.high\_wait\_time & 5.0 & s \\
    \hline
    vent2.dv\_max & 6.0 & m/s \\
    \hline
    \end{tabular}
\end{table}

\subsubsection{VentSet2}
The second vent set contains 6 TwoLevelVents built with the integrated constructor that all share the same properties as Vent1.

\newpage
\subsection{Nominal}
\textbf{RUN\_01\_nominal}

\noindent This test verifies all the main features of the two-level venting model. We randomize the two time parameters
 \emph{high\_duration} and \emph{wait\_time\_high} of Vent1. while keeping Vent2's constant. VentSet2 has six vents with
 non-randomized time parameters.

\subsubsection{Results}
This tests emphasizes that:
\begin{itemize} 
    \item the two separate vents Vent1 and Vent2 each produces high mode venting disturbances with unique durations and periods
    \item the behavior for each vent is independent of all the other vents
    \item the commanded duration adheres to the pre-defined minimum and maximum limits
    \item the magnitude of the venting force switches between a minimum and maximum value depending on whether it is in a high or low state
    \item the commanded wait time adheres to the pre-defined minimum and maximum limits
    \item the forces generated and collected are as expected
\end{itemize}

\begin{table}[!ht]
    \centering
    \begin{tabular}{|c|c|c|c|c|c|c|}
    \hline
    \textbf{Time (s)} & \textbf{Force[0] (N)} & \textbf{Force[1]} & \textbf{Force[2]} & \textbf{Torque[0](N.m)} & \textbf{Torque[1]} & \textbf{Torque[2]}\\
    \hline
    \end{tabular}
\end{table}

\begin{figure}[H]
\begin{center}
\includegraphics[scale=0.6]{Images/RUN_01_verif_set1.png}
\caption{Raw data for RUN\_01\_nominal - VentSet1}
\label{fig:RUN_01_verif_set1}
\end{center}
\end{figure}

\begin{table}[!ht]
    \centering
    \begin{tabular}{|c|c|c|c|c|c|c|}
    \hline
    \textbf{Time (s)} & \textbf{Force[0] (N)} & \textbf{Force[1]} & \textbf{Force[2]} & \textbf{Torque[0](N.m)} & \textbf{Torque[1]} & \textbf{Torque[2]}\\
    \hline
    \end{tabular}
\end{table}

\begin{figure}[H]
\begin{center}
\includegraphics[scale=0.6]{Images/RUN_01_verif_set2.png}
\caption{Raw data for RUN\_01\_nominal - VentSet2}
\label{fig:RUN_01_verif_set2}
\end{center}
\end{figure}

\begin{table}[!ht]
    \centering
    \resizebox{\textwidth}{!}{%
    \begin{tabular}{|c|c|c|c|c|c|c|c|}
    \hline
    \textbf{Time (s)} & \textbf{High venting} & \textbf{Force Mag (N)} & \textbf{High duration (s)} & \textbf{Wait time (s)} & \textbf{sched\_onTime (s)} & \textbf{DV mag (m/s)} & \textbf{start time high (s)}\\
    \hline
    \end{tabular}
    }
\end{table}

\begin{figure}[H]
\begin{center}
\includegraphics[scale=0.6]{Images/RUN_01_verif_vent1.png}
\caption{Raw data for RUN\_01\_nominal - Vent1}
\label{fig:RUN_01_verif_vent1}
\end{center}
\end{figure}

\begin{table}[!ht]
    \centering
    \resizebox{\textwidth}{!}{%
    \begin{tabular}{|c|c|c|c|c|c|c|c|}
    \hline
    \textbf{Time (s)} & \textbf{High venting} & \textbf{Force Mag (N)} & \textbf{High duration (s)} & \textbf{Wait time (s)} & \textbf{sched\_onTime (s)} & \textbf{DV mag (m/s)} & \textbf{start time high (s)}\\
    \hline
    \end{tabular}
    }
\end{table}

\begin{figure}[H]
\begin{center}
\includegraphics[scale=0.6]{Images/RUN_01_verif_vent2.png}
\caption{Raw data for RUN\_01\_nominal - Vent2}
\label{fig:RUN_01_verif_vent2}
\end{center}
\end{figure}




\newpage
\subsection{Fixed High Mode Duration}

\noindent \textbf{RUN\_02\_fixedDur}

\noindent This test sets a high mode duration to a fixed value and randomize the high mode wait time.
The changes from the default parameter table are shown below:

\begin{table}[!h]
    \centering
    \begin{tabular}{|c|c|c|} \hline
    \textbf{Parameter Name} & \textbf{Parameter Setting} & \textbf{Units}\\
    \hline
    wait\_time\_randFlag & TRUE & bool  \\
    \hline
    high\_duration & 2.0 & s  \\
    \hline
    wait\_time\_max & 5.0 & s  \\
    \hline
    wait\_time\_min & 3.0 & s  \\
    \hline
    dv\_max & 4.0 & m/s \\
    \hline
    \end{tabular}
\end{table}

\subsubsection{Results}
We cannot predict the random value high duration will have. Therefore, we simply provide here the raw data
that is simple enough to verify the model capabilities.

\begin{table}[!ht]
    \centering
    \begin{tabular}{|c|c|c|c|}
    \hline
    \textbf{Time (s)} & \textbf{high venting} & \textbf{high duration (s)} & \textbf{sched\_onTime (s) } \\
    \hline
    \end{tabular}
\end{table}

\begin{figure}[H]
\begin{center}
\includegraphics[scale=0.7]{Images/RUN_02_verif.png}
\caption{Raw data for RUN\_02\_fixedDur}
\label{fig:RUN_02_verif}
\end{center}
\end{figure}


\newpage
\subsection{Fixed High Mode Period}
\noindent \textbf{RUN\_03\_fixedWaitTime}

\noindent This test sets the high mode wait time to a fixed value while the duraiton is randomized. 
The changes from the default parameter table are shown below:

\begin{table}[!h]
    \centering
    \begin{tabular}{|c|c|c|} \hline
    \textbf{Parameter Name} & \textbf{Parameter Setting} & \textbf{Units}\\
    \hline
    high\_duration\_randFlag & TRUE & bool  \\
    \hline
    high\_wait\_time & 4.0 & s  \\
    \hline
    high\_duration\_max & 3.0 & s  \\
    \hline
    high\_duration\_min & 1.0 & s  \\
    \hline
    high\_wait\_time & 4.0 & s  \\
    \hline
    dv\_max & 6.0 & m/s \\
    \hline
    \end{tabular}
\end{table}

\subsubsection{Results}
We cannot predict the random value high duration will have. Therefore, we simply provide here the raw data
that is simple enough to verify the model capabilities.

\begin{table}[!ht]
    \centering
    \begin{tabular}{|c|c|c|c|c|}
    \hline
    \textbf{Time (s)} & \textbf{high venting} & \textbf{high duration (s)} & \textbf{DV mag (m/s)} & \textbf{sched\_onTime (s) } \\
    \hline
    \end{tabular}
\end{table}

\begin{figure}[H]
\begin{center}
\includegraphics[scale=0.7]{Images/RUN_03_verif.png}
\caption{Raw data for RUN\_03\_fixedWaitTime}
\label{fig:RUN_03_verif}
\end{center}
\end{figure}

\newpage
\subsection{Reset Event}
\textbf{RUN\_04\_reset\_event\_action}

\noindent This test verifies that the vent is correctly reset when the TwoLevelVentReset event-action
\emph{specific\_execution()} method is called. Duration and wait time are not randomized.

\subsubsection{Results}
    the \emph{sched\_onTime} is set to a different value. The value is not zero (as one could expect) because
    the value is logged at the end of every cycle, while the reset is done at the beginning. Therefore, the \emph{sched\_onTime}
    is updated before being logged. This behavior is desired.

\begin{table}[!ht]
    \centering
    \begin{tabular}{|c|c|c|c|c|}
    \hline
    \textbf{Time (s)} & \textbf{sched\_onTime (s)} & \textbf{start time high (s)} & \textbf{high venting} & \textbf{Force mag (N)} \\
    \hline
    0.0 & 5.0 & 0.0 & 1 & 200.0\\
    \hline
    1.0 & 5.0 & 0.0 & 1 & 200.0\\
    \hline
    2.0 & 5.0 & 0.0 & 0 & 100.0\\
    \hline
    3.0 & 5.0 & 0.0 & 0 & 100.0\\
    \hline
    4.0 & 5.0 & 0.0 & 0 & 100.0\\
    \hline
    5.0 & 10.0 & 5.0 & 1 & 200.0\\
    \hline
    6.0 & 10.0 & 5.0 & 1 & 200.0\\
    \hline
    7.0 & 10.0 & 5.0 & 0 & 100.0\\
    \hline
    8.0 & 13.0 & 8.0 & 1 & 200.0\\
    \hline
    9.0 & 13.0 & 8.0 & 1 & 200.0\\
    \hline
    10.0 & 13.0 & 8.0 & 0 & 100.0\\
    \hline
    \end{tabular}
\end{table}

\begin{figure}[H]
\begin{center}
\includegraphics[scale=0.6]{Images/RUN_04_verif.png}
\caption{Raw data for RUN\_04\_reset\_event\_action}
\label{fig:RUN_04_verif}
\end{center}
\end{figure}

\newpage
\subsection{Errors}
\textbf{ERROR\_rand\_config} \\
\textbf{ERROR\_rand\_duration} \\
\textbf{ERROR\_rand\_wait\_time} \\
\textbf{ERROR\_time\_config} \\

\noindent This section gathers runs that verify appropriate error messages are published when maximum high mode
duration is larger than minimum wait time or minimum high mode duration is larger for:
\begin{itemize}
    \item both wait time and duration randomized
    \item duration randomized only
    \item wait time randomized only
    \item no randomization
\end{itemize}

\subsection{Results}
\noindent All these cases output the same error message. Figure \ref{fig:unit_test_error_1} show the published error message. 

\begin{figure}[H]
\begin{center}
\includegraphics[scale=0.4]{Images/error_msg_1.png}
\caption{Error Message for invalid time inputs}
\label{fig:unit_test_error_1}
\end{center}
\end{figure}

\newpage
\subsection{Warnings}
\textbf{WARN\_01\_dur} \\
\textbf{WARN\_02\_wait\_time} \\
\textbf{WARN\_03\_maxDV} \\

\noindent This section gathers runs that verify appropriate warning messages are published when:
\begin{itemize}
    \item the maximum delta-V is reached, resetting the high mode duration to match maximum delta-V
    \item minimum high mode wait time is larger than the maximum high mode wait time when randomizing wait time,
        resetting the flag to False
    \item minimum high mode duration is larger than the maximum high mode duration when randomizing high duration,
        resetting the flag to False
\end{itemize}

\subsection{Results}
\noindent Figures \ref{fig:unit_test_warn_1}, \ref{fig:unit_test_warn_2} and \ref{fig:unit_test_warn_3} show the published error message. 

\begin{figure}[H]
\begin{center}
\includegraphics[scale=0.4]{Images/warn_msg_1.png}
\caption{Error Message for invalid randomized duration inputs}
\label{fig:unit_test_warn_1}
\end{center}
\end{figure}

\begin{figure}[H]
\begin{center}
\includegraphics[scale=0.4]{Images/warn_msg_2.png}
\caption{Error Message for invalid randomized wait time inputs}
\label{fig:unit_test_warn_2}
\end{center}
\end{figure}

\begin{figure}[H]
\begin{center}
\includegraphics[scale=0.4]{Images/warn_msg_3.png}
\caption{Warning Message for maximum delta-V reached}
\label{fig:unit_test_warn_3}
\end{center}
\end{figure}


\newpage
\subsection{Invalid Index Fail}

\noindent This test verifies that the run fails if the user tries to get a TwoLevelVent pointer at an incorrect index (that matches a SimpleVent).

\subsubsection{Results}
Figure \ref{fig:unit_test_fail} shows that:

\begin{itemize}
    \item the run is correctly failed.
\end{itemize}


\begin{figure}[H]
\begin{center}
\includegraphics[scale=0.3]{Images/fail.png}
\caption{Fail for invalid index in getter method}
\label{fig:unit_test_fail}
\end{center}
\end{figure}

\chapter{Future Work}

\noindent No future work.

\end{document}
